\chapter{Diseño del entorno de pruebas controlado}
\label{cap:diseno}

Este capítulo describe el diseño del entorno de pruebas (\breachlab)
creado con el fin de reproducir de forma controlada las seis vulnerabilidades
analizadas en el capítulo~\ref{cap:explotacion}, incluyendo en el proceso la justificación de las
decisiones de diseño,
la arquitectura resultante, el modelo de datos, el mecanismo que permite
alternar entre la rama vulnerable y la rama segura del código, el registro
de evidencias y el procedimiento de despliegue.

\section{Justificación del diseño}

Existen laboratorios de referencia ampliamente utilizados en formación en
ciberseguridad, como DVWA~\citelink{DVWA} u OWASP Juice
Shop~\citelink{JUICESHOP} (véase el capítulo~\ref{cap:comparativa}), que
cubren un catálogo de vulnerabilidades más amplio que el de este trabajo.
Sin embargo, ambos reproducen arquitecturas monolíticas en las que el
servidor genera el HTML final. Una parte creciente de las aplicaciones
reales sigue en cambio un patrón de interfaz de tipo SPA (\emph{Single Page
Application}) que consume una API JSON independiente, patrón que modifica de
forma sustancial cómo se manifiestan vulnerabilidades como el XSS (que pasa
a ser un problema del DOM en el cliente y no de una plantilla en el
servidor). \breachlab se ha diseñado expresamente con esta arquitectura para
poder discutir esa diferencia con conocimiento de causa.

Un segundo requisito de diseño fue servir el frontend y el backend bajo el
mismo origen. Esto es imprescindible para poder demostrar el CSRF de forma
realista: si la API se sirviera en un origen distinto sujeto a CORS, el
navegador ya estaría aplicando una restricción adicional que no reflejaría
el caso típico de una aplicación desplegada como monolito lógico (aunque sea
físicamente en dos contenedores).

\section{Arquitectura}

\breachlab está formado por tres componentes desplegados mediante Docker
Compose: un contenedor nginx que sirve los artefactos estáticos generados
por Astro; un contenedor
Flask~\citelink{FLASK} que expone la API JSON de la aplicación; y una base
de datos SQLite persistida
en un volumen. 

\needspace{15\baselineskip}
El navegador solo habla con nginx, a través del puerto 8080: las
peticiones a rutas estáticas (\verb|/|) se resuelven contra el build de
Astro y las peticiones a \verb|/api/*| se reenvían al backend Flask. Servir
ambos componentes bajo el mismo origen es lo que permite que la cookie de
sesión viaje sin restricciones de CORS entre frontend y backend.

\begin{figure}[htbp]
  \centering
  \begin{tikzpicture}[
      node distance=9mm and 14mm,
      box/.style={draw, rounded corners, minimum width=3.1cm,
        minimum height=9mm, align=center, font=\small},
      lbl/.style={font=\scriptsize, align=center, fill=white, inner sep=1pt},
      arr/.style={-{Latex[length=2mm]}}
    ]
    \node[box] (browser) {Navegador};
    \node[box, below=of browser] (nginx) {nginx (web)\\[1pt]\texttt{:8080}};
    \node[box, below left=16mm and -2mm of nginx] (astro) {Astro\\[1pt](\emph{build} estático)};
    \node[box, below right=16mm and -2mm of nginx] (flask) {Flask\\[1pt]\texttt{api:5000}};
    \node[box, below=of flask] (sqlite) {SQLite};

    \draw[arr] (browser) -- node[lbl, right] {HTTP\\(mismo origen)} (nginx);
    \draw[arr] (nginx.south) -- ++(0,-4mm) -| node[lbl, pos=0.6, above] {\texttt{/}\\(estático)} (astro.north);
    \draw[arr] (nginx.south) -- ++(0,-4mm) -| node[lbl, pos=0.6, above] {\texttt{/api/*}\\(proxy)} (flask.north);
    \draw[arr] (flask) -- (sqlite);
  \end{tikzpicture}
  \caption{Arquitectura de \breachlab.}
  \label{fig:arquitectura}
\end{figure}

El interruptor \texttt{?safe=1} es la decisión de diseño más interesante del
laboratorio: añadido a la URL de cualquier página, se propaga tanto a las
llamadas del frontend a la API como a la lógica de renderizado del propio
frontend, activando de forma simultánea la contramedida de backend y de
frontend para la vulnerabilidad correspondiente. Asimismo, el
código fuente contiene ambas ramas (vulnerable y corregida) de cada
endpoint, lo que permite comparar directamente el antes y el después sin
mantener dos aplicaciones independientes.

\section{Modelo de datos y usuarios de prueba}

Se ha optado por SQLite como motor de base de datos por su sencillez: no
requiere un servidor independiente ni configuración previa, y resulta
suficiente para un laboratorio de un único usuario concurrente con finalidad
educativa y no de rendimiento.

El esquema de la base de datos consta de dos tablas: \verb|users| (\verb|id|,
\verb|username|, \verb|password| en texto plano, \verb|password_hash| y
\verb|role|) y \verb|notes| (\verb|id|, \verb|owner_id|, \verb|title|,
\verb|content|). El campo \verb|password| se mantiene deliberadamente en
texto plano junto al hash para poder ilustrar en el
capítulo~\ref{cap:explotacion} la diferencia
entre la ruta vulnerable, que compara contraseñas en claro, y la ruta
segura, que utiliza \verb|check_password_hash| del módulo de seguridad de
Werkzeug~\citelink{WERKZEUG}.

\begin{table}[htbp]
  \centering
  \begin{tabular}{lll}
    \toprule
    \tabhead{Usuario} & \tabhead{Contraseña} & \tabhead{Rol} \\
    \midrule
    admin  & Admin123!     & admin \\
    marina & contraSegura  & user  \\
    pablo  & qwerty        & user  \\
    \bottomrule
  \end{tabular}
  \caption{Usuarios de prueba precargados por \texttt{start\_db()} en \texttt{backend/api.py}.}
  \label{tab:usuarios}
\end{table}

Además de los usuarios, se precargan dentro de la base de datos cuatro notas con datos ficticios
 (una tarjeta de crédito, la contraseña de un correo, etc.),
propiedad de distintos usuarios, que se utilizan en el
capítulo~\ref{cap:explotacion} para demostrar el acceso indebido a
información de terceros a través del IDOR.

\section{Registro de evidencias}

Cada acción relevante (inicio de sesión correcto o fallido, inyección SQL,
cambio de contraseña, acceso a notas ajenas, acceso al panel de
administración) permanece registrado mediante la función \verb|log_event()| de
\texttt{backend/api.py} en el fichero \texttt{backend/logs/attacks.log},
persistido fuera del contenedor mediante un volumen Docker. Cada línea
incluye una marca de tiempo, el modo de ejecución (vulnerable o seguro), la
IP de origen, el tipo de evento y detalles específicos del mismo (usuario,
identificador de nota, consulta SQL ejecutada, etc.). Este registro se
utiliza en el capítulo~\ref{cap:explotacion} como evidencia objetiva de cada
explotación.

\section{Despliegue}

El laboratorio se levanta con un único comando, \verb|docker compose up --build|,
quedando disponible en la máquina local, accediendo desde \url{http://localhost:8080}. 
Otra opción sería ejecutar el backend Flask y el frontend Astro por separado, en
cuyo caso el proxy configurado en \texttt{astro.config.mjs} redirige las
llamadas a \texttt{/api} hacia el backend en el puerto 5000. 

El endpoint \verb|GET /api/reset| reinicializa la base de datos a su estado inicial, lo
que resulta útil para repetir una explotación en condiciones limpias antes
de capturar evidencias.
